% Original vector schematic of the FLFM light path (not to scale).
% Symbols follow Section "Abbreviated Calculation Section" of Chapter 3.
\begin{tikzpicture}[
    x=1cm,y=1cm,
    font=\footnotesize,
    ink/.style={text=black!85},
    lens/.style={<->,>={Stealth[length=2mm]},draw=black!75,line width=0.8pt},
    ray/.style={draw=blue!60!black,line width=0.7pt},
    slice/.style={draw=black!25,line width=0.4pt},
    plane/.style={draw=black!45,dashed,line width=0.5pt},
    dim/.style={<->,>={Stealth[length=1.6mm]},draw=black!70,line width=0.5pt},
    lab/.style={font=\footnotesize,ink,inner sep=1.5pt,align=center}
]
% ---- geometry (cm; not to scale) ----
% object 0 | objective 1.3 | pupil 2.6 | tube lens 4.6 | NIP 6.6
% | Fourier lens 9.9 | MLA 13.2 | camera plane 14.3
\def\xo{1.3}   \def\xp{2.6}   \def\xt{4.6}  \def\xn{6.6}
\def\xf{9.9}   \def\xm{13.2}  \def\xc{14.3}
\def\hp{0.8}   % half of D_pupil
\def\hq{1.32}  % half of D_proj
\def\dm{0.528} % lenslet pitch d_MLA (5 lenslets across D_proj)

% optical axis
\draw[draw=black!30,line width=0.4pt] (-0.4,0) -- (\xc+0.4,0);

% planes
\draw[plane] (0,-1.5) -- (0,1.5);
\draw[plane] (\xp,-1.5) -- (\xp,1.5);
\draw[plane] (\xn,-1.5) -- (\xn,1.5);
\draw[plane] (\xc,-1.5) -- (\xc,1.5);

% lenses
\draw[lens] (\xo,-1.15) -- (\xo,1.15);
\draw[lens] (\xt,-1.15) -- (\xt,1.15);
\draw[lens] (\xf,-1.6) -- (\xf,1.6);

% microlens array: 7 lenslets, 5 of them filled by the pupil image
\foreach \k in {-3,...,3} {
    \draw[draw=black!70,fill=black!6,line width=0.5pt]
        (\xm-0.07,{\k*\dm-0.5*\dm}) rectangle (\xm+0.07,{\k*\dm+0.5*\dm});
}

% rays of one on-axis emitter (marginal rays)
\draw[ray] (0,0) -- (\xo,\hp);
\draw[ray] (0,0) -- (\xo,-\hp);
\draw[ray] (\xo,\hp) -- (\xt,\hp);
\draw[ray] (\xo,-\hp) -- (\xt,-\hp);
\draw[ray] (\xt,\hp) -- (\xn,0);
\draw[ray] (\xt,-\hp) -- (\xn,0);
\draw[ray] (\xn,0) -- (\xf,\hq);
\draw[ray] (\xn,0) -- (\xf,-\hq);
\draw[ray] (\xf,\hq) -- (\xm,\hq);
\draw[ray] (\xf,-\hq) -- (\xm,-\hq);

% slicing of the collimated beam into sub-beams, one per lenslet,
% each focused to a diffraction-limited spot on the camera plane
\foreach \k in {-2,...,2} {
    \draw[slice] (\xf,{\k*\dm-0.5*\dm}) -- (\xm,{\k*\dm-0.5*\dm});
    \draw[ray] (\xm,{\k*\dm-0.5*\dm}) -- (\xc,{\k*\dm});
    \draw[ray] (\xm,{\k*\dm+0.5*\dm}) -- (\xc,{\k*\dm});
    \fill[orange!80!black] (\xc,{\k*\dm}) circle (0.045);
}
\draw[slice] (\xf,{2.5*\dm}) -- (\xm,{2.5*\dm});

% dimensions
\draw[dim] (\xp,-\hp) -- (\xp,\hp);
\node[lab,anchor=east,fill=white,inner sep=1pt] at (\xp-0.05,-0.35) {$D_{pupil}$};
\draw[dim] (11.4,-\hq) -- (11.4,\hq);
\node[lab,anchor=east,fill=white,inner sep=1pt] at (11.35,0.0) {$D_{proj}$};
\draw[dim] (\xm-0.45,{1.5*\dm}) -- (\xm-0.45,{2.5*\dm});
\node[lab,anchor=east,fill=white,inner sep=1pt] at (\xm-0.5,{2.0*\dm}) {$d_{MLA}$};

% focal lengths along the bottom
\def\yb{-2.15}
\draw[dim] (0,\yb) -- (\xo,\yb);   \node[lab,fill=white] at (0.65,\yb) {$f_{obj}$};
\draw[dim] (\xt,\yb) -- (\xn,\yb); \node[lab,fill=white] at (5.6,\yb)  {$f_{TL}$};
\draw[dim] (\xn,\yb) -- (\xf,\yb); \node[lab,fill=white] at (8.25,\yb) {$f_{FL}$};
\draw[dim] (\xf,\yb) -- (\xm,\yb); \node[lab,fill=white] at (11.55,\yb) {$f_{FL}$};
\draw[dim] (\xm,\yb) -- (\xc,\yb); \node[lab,fill=white] at (13.75,\yb) {$f_{MLA}$};

% names
\node[lab,anchor=south] at (0,1.55) {object\\plane};
\node[lab,anchor=south] at (\xo,1.2) {objective};
\node[lab,anchor=south] at (\xp+0.1,1.55) {back\\pupil};
\node[lab,anchor=south] at (\xt,1.2) {tube lens\\(TL)};
\node[lab,anchor=south] at (\xn,1.55) {NIP};
\node[lab,anchor=south] at (\xf,1.65) {Fourier lens\\(FL)};
\node[lab,anchor=south] at (\xm-0.35,2.05) {MLA};
\node[lab,anchor=south west] at (\xc-0.1,1.55) {camera plane\\(MLA focal\\plane)};

% spot size
\node[lab,anchor=west,align=left] at (\xc+0.35,-0.9)
    {spot\\[-1pt]$R_{xy}^{MLA}=\dfrac{\lambda}{2\,NA_{MLA}}$};
\draw[draw=orange!80!black,line width=0.4pt] (\xc+0.33,-0.7) -- (\xc+0.06,{-2*\dm-0.02});

% pupil-sampling note
\node[lab,anchor=north] at (12.3,-1.5)
    {$N = D_{proj}/d_{MLA}$ lenslets across the pupil image};
\end{tikzpicture}
